% ============================================================
% 5. DISCUSSION -- humanized. Cut self-congratulation ("worth naming",
% "we think one of the transferable lessons"). GABRIL: state direction,
% drop the second-hand-ness advertisement (fix #4).
% ============================================================
\section{Discussion}

\paragraph{What the null does and does not say.}
The central result is a null with a mechanism behind it. Conditioning a frozen
V-JEPA~2 predictor on gaze and hand at the input does not improve short-horizon
anticipation, and the five diagnostics say why: the signal reaches the
predictor's output, the predictor does not respond to its content, and on its own
the signal predicts the next action barely above a class prior. So the claim we
are entitled to is narrow. Not that gaze and hand are uninformative, but that at
this horizon, this sample size, and through this pathway they carry too little
action-relevant information to move the prediction. That narrower claim matters,
because it points to a different fix than ``build a bigger pathway.''

\paragraph{Detectability versus decision-relevance.}
A single accuracy number hides a distinction the PEVA-style readout makes
visible. One question is whether the signal is detectable at all; another is
whether it is large enough to change a decision. Here the signal is detectable in
the statistical sense (it beats the masked version on 172 of 220 clips) yet
produces no discriminative margin (not one discordant pair). An effect that is
real but decision-irrelevant is easy to report as a success if only average error
is measured, which is one reason we ran both stages of the readout rather than
the first alone.

\paragraph{The motion/object dissociation, in context.}
The one piece of residual structure is that gaze and hand favour verbs while
vision favours nouns. The natural reading is that behavioral signals encode how
the body is about to move rather than what it will act on, and that EK100 is hard
mainly because of object identity. As a performance result this carries little
weight, since it sits below the verb prior. As corroboration it carries more: the
two strongest published EK100 systems independently split their probes along the
same verb/noun line and defend it with the same motion-versus-appearance argument
\citep{chu2026jfaa, wang2026tapjepa}. When three separate efforts land on the
same boundary, it is likely a property of the task rather than of our setup.

\paragraph{Why input-level conditioning may be the weak use of these signals.}
Set against the recent literature, our null looks less like a surprise and more
like one more instance of a pattern. Methods that use gaze to \emph{supervise} a
representation during training and then drop it at inference tend to report gains
where input concatenation, as here, does not: the gaze regulariser of
GABRIL \citep{banayeeanzade2025gabril}, the gaze-regularised vision-language
models of \citet{pani2026gazereg}, and the goal-consistency objective of
\citet{royfernando2022abstractgoal} all take this route. GazeQwen puts the point
directly, arguing that the obstacle is a missing learned pathway from gaze into
the model's internal representations rather than any deficit in the signal. Our
diagnostics fit that argument at the input level almost exactly: the model
attends to where the signal sits but not to what it says, which is what a missing
learned pathway would look like.

\paragraph{Relation to the horizon analysis.}
The horizon sweep of \S\ref{subsec:horizon} measures something different from the
margins above: how far ahead a behavioral target is recoverable from the frozen
representation at all. It finds that behavioral legibility is a present-tense
property, largely gone by one to two seconds, which bounds what any input-level
pathway could have delivered at the horizon we probe. One caution keeps the two
analyses consistent. The behavior-minus-vision margin in \S\ref{subsec:dissociation} does widen with
horizon, but the widening splits roughly evenly between a decaying vision
comparator and a rising behavioral arm, and the behavioral arm remains below its
own class prior at every horizon. Since neither published challenge system sweeps
horizon, there is no outside evidence for a ``signal grows with horizon'' reading
either. The
horizon numbers describe the regime; they are not evidence that a longer horizon
rescues the signal.